% !TEX root = ../algebra_02.tex

\section{Подполя поля из $p^n$ элементов}

\begin{Theorem}{Структура подполей конечного поля}{}
	Пусть $p$ --- простое число, $m, n \in \mathbb{N}$.
	\begin{enumerate}
		\item Если $m \nmid n$, то в поле $\mathbb{F}_{p^n}$ нет подполей из $p^m$ элементов.
		\item Если $m \mid n$, то в поле $\mathbb{F}_{p^n}$ существует единственное подполе из $p^m$ элементов.
	\end{enumerate}
\end{Theorem}

\begin{proof}
	\begin{enumerate}
		\item Пусть $K \subset \mathbb{F}_{p^n}$ --- подполе, $|K| = p^m$. Тогда $\mathbb{F}_{p^n}$ можно рассматривать как линейное пространство над полем $K$. Пусть его размерность равна $[\mathbb{F}_{p^n} : K] = d \in \mathbb{N}$. Тогда:
		\[ |\mathbb{F}_{p^n}| = |K|^d \implies p^n = (p^m)^d = p^{md} \]
		Следовательно, $n = md$, откуда $m \mid n$. Противоречие.

		\item Докажем единственность. Пусть $K \subset \mathbb{F}_{p^n}$ --- подполе из $p^m$ элементов. Тогда для любого $a \in K$ выполнено $a^{p^m} - a = 0$. Следовательно, $K$ полностью содержится во множестве корней многочлена $X^{p^m}-X$:
		\[ K \subset \{ a \in \mathbb{F}_{p^n} \mid a^{p^m} - a = 0 \} = F \]
		Поскольку многочлен степени $p^m$ имеет не более $p^m$ корней, а множество $K$ содержит ровно $p^m$ элементов, то $K = F$. Значит, если подполе существует, оно состоит ровно из корней этого многочлена и потому единственно.

		Само множество $F$ образует поле (доказывается аналогично построению поля разложения для эндоморфизма Фробениуса), что доказывает и существование.
	\end{enumerate}
\end{proof}
